%==============================================================================
% Sjabloon onderzoeksvoorstel bachproef
%==============================================================================
% Gebaseerd op document class `hogent-article'
% zie <https://github.com/HoGentTIN/latex-hogent-article>

% Voor een voorstel in het Engels: voeg de documentclass-optie [english] toe.
% Let op: kan enkel na toestemming van de bachelorproefcoördinator!
\documentclass{hogent-article}

% Invoegen bibliografiebestand
\addbibresource{voorstel.bib}

% Informatie over de opleiding, het vak en soort opdracht
\studyprogramme{Professionele bachelor toegepaste informatica}
\course{Bachelorproef}
\assignmenttype{Onderzoeksvoorstel}
% Voor een voorstel in het Engels, haal de volgende 3 regels uit commentaar
% \studyprogramme{Bachelor of applied information technology}
% \course{Bachelor thesis}
% \assignmenttype{Research proposal}

\academicyear{2025-2026} % TODO: pas het academiejaar aan

% TODO: Werktitel
\title{Detectie van AI-gebaseerde Cheats in Videogames via Gedragsanalyse: Een Case Study in Garry’s Mod}

% TODO: Studentnaam en emailadres invullen
\author{Gauthier Vandeputte}
\email{Gauthier.Vandeputte@student.hogent.be}

% TODO: Medestudent
% Gaat het om een bachelorproef in samenwerking met een student in een andere
% opleiding? Geef dan de naam en emailadres hier
% \author{Yasmine Alaoui (naam opleiding)}
% \email{yasmine.alaoui@student.hogent.be}

% TODO: Geef de co-promotor op
\supervisor[Co-promotor]{}

% Binnen welke specialisatierichting uit 3TI situeert dit onderzoek zich?
% Kies uit deze lijst:
%
% - Mobile \& Enterprise development
% - AI \& Data Engineering
% - Functional \& Business Analysis
% - System \& Network Administrator
% - Mainframe Expert
% - Als het onderzoek niet past binnen een van deze domeinen specifieer je deze
%   zelf
%
\specialisation{AI \& Data Engineering}
\keywords{AI, Deep Learning, cheats, videogames}

\begin{document}

\begin{abstract}
Het doel van dit onderzoek is het verkrijgen van diepgaand inzicht in de ontwikkeling van AI-gebaseerde cheats en het identificeren van effectieve methoden om deze tegen te gaan. In dit kader is de volgende onderzoeksvraag geformuleerd: Hoe worden AI-gebaseerde cheats ontwikkeld en welke methoden kunnen worden ingezet om deze te detecteren op basis van spelersgedrag en interactiepatronen? In dit onderzoek wordt onder spelersgedrag verstaan: de fysieke muis- en toetsenbordbewegingen van spelers, evenals de reactiesnelheid waarmee zij op specifieke spelsituaties reageren.
\\\\
Om de onderzoeksvraag te kunnen beantwoorden, wordt eerst onderzocht op welke wijze muis- en toetsenbordinvoer kan worden vervalst met behulp van een tweede computer. Vervolgens wordt geanalyseerd hoe een model kan worden getraind voor het detecteren en interpreteren van relevante spelinstanties en het hierop adequaat reageren. Tot slot wordt een analyse uitgevoerd op de verzamelde gebruikersinvoer om anomalieën te identificeren die kunnen wijzen op het gebruik van AI-gebaseerde cheats.
\\\\
De resultaten van dit onderzoek zullen leiden tot de ontwikkeling van een raamwerk dat de specifieke gedrags- en inputkenmerken identificeert die duiden op het gebruik van AI-gebaseerde cheats. Op basis hiervan zal een onderbouwde aanbeveling worden gedaan voor de meest geschikte machine learning-methoden om afwijkend spelersgedrag te detecteren, waarmee een effectieve aanpak voor het bestrijden van dit probleem kan worden gerealiseerd.
\end{abstract}

\tableofcontents

% De hoofdtekst van het voorstel zit in een apart bestand, zodat het makkelijk
% kan opgenomen worden in de bijlagen van de bachelorproef zelf.
%---------- Inleiding ---------------------------------------------------------

\section{Inleiding}%
\label{sec:inleiding}

In de wereld van competitieve games is een eerlijk speelveld essentieel. Dit vormt het fundament van de hedendaagse gamingindustrie, waarin miljarden euro’s omgaan. De afgelopen jaren is er echter een zorgwekkende verschuiving waarneembaar: de opkomst van valsspelen via Artificiële Intelligentie (AI). Dit vormt een aanzienlijk probleem, aangezien traditionele anti-cheat-systemen die speuren naar gemodificeerde gamebestanden niet in staat zijn deze nieuwe vorm van fraude te detecteren.
\\\\
Deze AI-gebaseerde cheats opereren op een geavanceerde wijze: er wordt veelal gebruikgemaakt van een externe computer om de beelden op het scherm te analyseren, waarna menselijke input (zoals muisbewegingen en toetsaanslagen) wordt gesimuleerd. Doordat de cheat onafhankelijk van de game-software functioneert, blijft deze nagenoeg onzichtbaar voor gangbare beveiliging. De doelgroep voor dit onderzoek bestaat uit IT-professionals en specialisten op het gebied van anti-cheat-software binnen de gamingindustrie. Dit toegepaste onderzoek beoogt een concrete oplossing te bieden voor dit groeiende probleem.
\\\\
Hiertoe analyseert dit onderzoek het gedrag van de speler. Data met betrekking tot muisbewegingen en reactiesnelheid worden onderzocht om patronen te identificeren die afwijken van hetgeen fysiek realistisch is voor een menselijke gebruiker. Dit leidt tot de centrale onderzoeksvraag:
\\\\
\textit{Hoe worden AI-gebaseerde cheats ontwikkeld en welke methoden, gebaseerd op de analyse van spelersgedrag en interactiepatronen, kunnen worden aangewend om deze effectief te identificeren?}
\\\\
Het uiteindelijke doel is het leveren van een bruikbaar resultaat in de vorm van een rapport met onderbouwde aanbevelingen. Hiermee wordt aangetoond welke Machine Learning-technieken het meest geschikt zijn voor het herkennen van afwijkend, geautomatiseerd spelersgedrag, om zodoende bij te dragen aan een eerlijkere online spelomgeving.

%---------- Stand van zaken ---------------------------------------------------

\section{Literatuurstudie}
\label{sec:literatuurstudie}

De opkomst van AI-gebaseerde cheats is een van de meest uitdagende ontwikkelingen binnen het domein van anti-cheat-beveiliging. Conventionele oplossingen richten zich primair op de detectie van gemodificeerde bestanden, verdachte processen of manipulaties in het computergeheugen. Deze methoden schieten echter tekort wanneer de cheat buiten de directe game-omgeving opereert. De zogenaamde "external cheats" functioneren volledig onafhankelijk van de spelsoftware en communiceren uitsluitend via beeldherkenning en gesimuleerde gebruikersinput, wat detectie door traditionele mechanismen bemoeilijkt \autocite{Kunstner2024}.
\\\\
Een significante verschuiving wordt veroorzaakt door de integratie van Machine Learning-technieken in moderne cheats. Onderzoek toont aan dat neurale netwerken niet alleen in staat zijn om tegenstanders in realtime te detecteren, maar ook om autonome muisbewegingen te genereren \autocite{Neuman2023}. Deze systemen verwerken live beelden via objectdetectie-algoritmen (zoals YOLO of EfficientDet) en vertalen de resultaten direct naar acties, zoals het automatisch richten op tegenstanders of het compenseren van de 'recoil' (terugslag). Doordat deze systemen zijn geoptimaliseerd voor een lage latentie en hoge precisie, kunnen zij menselijke interactie nabootsen op een wijze die handmatig nauwelijks te onderscheiden is.
\\\\
Terwijl interne cheats doorgaans sporen achterlaten via technieken zoals DLL-injectie, maken externe cheats gebruik van hardware of een tweede computer om de beelduitvoer te analyseren (bijvoorbeeld via HDMI-capture). Vervolgens wordt gesimuleerde input via HID-emulatie naar de primaire computer verzonden \autocite{Liu2020}. Omdat deze invoer voor het besturingssysteem identiek lijkt aan legitieme menselijke handelingen, vormt dit een fundamenteel probleem voor bestaande architecturen.
\\\\
Onderzoek naar de detectie van dergelijk gedrag richt zich daarom in toenemende mate op gedragsanalyse. Pinto et al. stellen dat tijdreeksmodellen uitermate geschikt zijn voor het onderscheiden van menselijke versus geautomatiseerde input, op basis van kenmerken zoals reactiesnelheid en de kinetiek van muisbewegingen \autocite{Pinto2021}. Deep learning-modellen, specifiek convolutionele en recurrente netwerken, behalen hierbij een significant hogere nauwkeurigheid dan statistische methoden.
\\\\
Literatuur over menselijke motoriek bevestigt dat natuurlijke bewegingen nooit volledig lineair of perfect reproduceerbaar zijn; ze vertonen variaties in acceleratie en timing die complex zijn om te simuleren \autocite{VanderKemp2019}. Wanneer AI-cheats deze patronen trachten te imiteren, ontstaan vaak subtiele inconsistenties, zoals onnatuurlijk vloeiende trajecten of identieke sequenties in opeenvolgende acties.
\\\\
Door de inzet van anomaliedetectie-algoritmen, zoals Isolation Forests of auto-encoders, kunnen afwijkingen worden geïdentificeerd zonder dat een specifieke cheat-strategie vooraf bekend hoeft te zijn \autocite{Chen2022}. Dit is essentieel in een snel evoluerend domein. Ondanks bestaand onderzoek naar cheatdetectie is er momenteel een gebrek aan literatuur die zich specifiek richt op externe, AI-gebaseerde cheats die steunen op beeldherkenning. Daarnaast ontbreekt een systematische evaluatie van het gedrag van dergelijke cheats binnen specifieke game-engines, wat de relevantie van dit onderzoek onderbouwt.

%---------- Methodologie ------------------------------------------------------
\section{Methodologie}%
\label{sec:methodologie}

Het onderzoek combineert de ontwikkeling van AI-gebaseerde cheats met de daaropvolgende detectie ervan. De focus ligt op technische diepgang en meetbare resultaten. De totale looptijd van dit onderzoek wordt geschat op twaalf weken.

\subsection{1. Opbouw van de projectomgeving}
\begin{itemize}
    \item \textbf{Doel:} Onderzoeken hoe computerinput (muis en toetsenbord) gesimuleerd kan worden met behulp van een tweede computer.
    \item \textbf{Stappen:}
    \begin{enumerate}
        \item \textbf{Hardware-analyse:} Onderzoek naar de verschillende hardwarematige mogelijkheden om gebruikersinput te emuleren.
        \item \textbf{Ontwikkeling testomgeving:} Het opzetten van een omgeving waarin beeldmateriaal wordt uitgelezen, het cheat-model wordt getraind en gebruikersinput wordt gelogd voor analyse.
    \end{enumerate}
    \item \textbf{Duur en Deliverables:}
    \begin{itemize}
        \item \textbf{Duur:} 3 weken.
        \item \textbf{Deliverable:} Een functionele opstelling voor beeldanalyse op een tweede computer, inclusief een testomgeving met logfunctionaliteit.
    \end{itemize}
    \item \textbf{Tools:} Raspberry Pi, VS Code.
\end{itemize}

\subsection{2. Ontwikkelen van AI-gebaseerde cheats}
\begin{itemize}
    \item \textbf{Doel:} Het ontwikkelen van de benodigde scripts voor de AI-gebaseerde cheat-functionaliteiten.
    \item \textbf{Stappen:}
    \begin{enumerate}
        \item \textbf{Recoil-compensatiescript:} Ontwikkeling van een script dat de terugslagpatronen van wapens neutraliseert door tegenovergestelde bewegingen te genereren.
        \item \textbf{Objectdetectiemodel:} Het trainen van een model voor de detectie van tegenstanders in diverse omgevingen, inclusief situaties waarbij de vijand deels achter obstakels verscholen gaat.
        \item \textbf{Realtime reactiescript:} Een script dat de live-beelden analyseert en visuele indicatoren (bounding boxes) plaatst bij gedetecteerde entiteiten.
    \end{enumerate}
    \item \textbf{Duur en Deliverables:}
    \begin{itemize}
        \item \textbf{Duur:} 3 weken.
        \item \textbf{Deliverable:} Drie gevalideerde scripts voor recoil-controle en vijanddetectie.
    \end{itemize}
    \item \textbf{Tools:} VS Code.
\end{itemize}

\subsection{3. Integratie en datalogging}
\begin{itemize}
    \item \textbf{Doel:} Het samenvoegen van de afzonderlijke scripts tot één coherent systeem en het verzamelen van logdata.
    \item \textbf{Stappen:}
    \begin{enumerate}
        \item \textbf{Systeemintegratie:} Koppeling van de scripts zodat het systeem autonoom kan richten en schieten zodra een doelwit wordt geïdentificeerd, waarbij de gebruiker de controle herwint na eliminatie van het doelwit.
        \item \textbf{Data-acquisitie:} Het loggen van een combinatie van menselijke en AI-gebaseerde input (muisbewegingen, reactietijden) om een dataset te genereren voor het detectiemodel.
    \end{enumerate}
    \item \textbf{Duur en Deliverables:}
    \begin{itemize}
        \item \textbf{Duur:} 3 weken.
        \item \textbf{Deliverable:} Een geïntegreerd softwaresysteem en een dataset in CSV-formaat.
    \end{itemize}
    \item \textbf{Tools:} VS Code.
\end{itemize}

\subsection{4. Model-evaluatie en rapportage}
\begin{itemize}
    \item \textbf{Doel:} Het trainen van een anti-cheat-model en het uitvoeren van analyses op de verzamelde data.
    \item \textbf{Stappen:}
    \begin{enumerate}
        \item \textbf{Training detectiemodel:} Gebruik van de dataset om een model te trainen dat anomalieën in inputpatronen herkent.
        \item \textbf{Realtime validatie:} Het model live testen tijdens gameplay om de effectiviteit van de anomaliedetectie te verifiëren.
    \end{enumerate}
    \item \textbf{Duur en Deliverables:}
    \begin{itemize}
        \item \textbf{Duur:} 3 weken.
        \item \textbf{Deliverable:} Een werkend prototype van een anti-cheat-systeem en een analyserapport.
    \end{itemize}
    \item \textbf{Tools:} VS Code.
\end{itemize}

%---------- Verwachte resultaten ----------------------------------------------
\section{Verwacht resultaat en conclusie}%
\label{sec:verwachte_resultaten}

De verwachting is dat dit onderzoek resulteert in een functionele AI-gebaseerde cheat die voldoet aan de technische specificaties en bovendien flexibel inzetbaar is voor verschillende speltitels. Tevens wordt verwacht dat de verzamelde data leiden tot een robuust anti-cheat-model dat in staat is om geautomatiseerde interacties met hoge accuraatheid te onderscheiden van menselijk handelen.

\printbibliography[heading=bibintoc]

\end{document}